\documentclass[12pt]{article}
\usepackage[utf8]{inputenc}
\usepackage[T1]{fontenc}
\usepackage{hyperref}
\title{Yalaku Language Resources for K-12 Curriculum Development}
\author{Compiled by Claude Sonnet 5.5 (claude-sonnet-5-5) for Andr\'as Kornai}
\date{}
\begin{document}
\maketitle

\section{Introduction}
Yalaku (ISO 639-3: ylg; formerly Yelogu, called Kaunga by the Kwoma) is a Ndu language of the Middle Sepik, East Sepik Province, Papua New Guinea. Its speakers live in three villages off the Sepik River: Yalaku, Hambakaini and Kumajuwi. Speaker estimates vary widely between sources, from about 300 to about 1,000. Yalaku has long been in contact with the unrelated Kwoma language, and Tok Pisin is widely known. The main recent research is by Alexandra Aikhenvald.

\section{Primary Educational Resources}
No primers, textbooks or school materials in Yalaku turned up in the search.

\section{Linguistic and Reference Materials}
\begin{itemize}
\item Aikhenvald, ``Comparison, contrast and similarity in Yalaku'': \url{https://researchonline.jcu.edu.au/54532/}
\item Aikhenvald, ``Language contact and language change in the Sepik Region of New Guinea: the case of Yalaku'': \url{https://researchonline.jcu.edu.au/63889/}
\item ``Word in Yalaku'' (book chapter): \url{https://acquire.cqu.edu.au/articles/chapter/Word_in_Yalaku/16780645}
\item \emph{Linguistic Discovery} article (2018) on Yalaku: \url{https://doaj.org/article/68115a2c97454c6a9f15305a44e920be}
\item Numerals entry: \url{https://lingweb.eva.mpg.de/channumerals/Yalaku.htm}
\item ASJP word-list entry: \url{https://asjp.clld.org/languages/YELOGU}
\end{itemize}
The search found no dictionary or pedagogical grammar. Aikhenvald notes that Laycock's 1965 outline contains many mistakes.

\section{Religious and Community Resources}
\begin{itemize}
\item Oceanic Art Society, ``The Yalaku: history and warfare in the Middle Sepik'': \url{https://www.oceanicartsociety.org.au/the-yalaku-history-and-warfare-in-the-middle-sepik/}
\item Endangered Languages Project entry: \url{https://endangeredlanguages.com/lang/10615?hl=fr}
\end{itemize}
No scripture resources turned up in the search.

\section{Digital and Audio Resources}
No audio, video or software resources turned up in the search.

\section{Curriculum Development Context}
The scholarly descriptions give a base for developing an orthography and teacher materials. The community is small and multilingual (Yalaku, Kwoma, Tok Pisin), so any programme would have to take that into account. No ready-made school materials were found.

\end{document}